% Spin(10) axion v20 — decay-safe anomaly completion and exact threshold matching
\documentclass[11pt,a4paper]{article}
\usepackage[margin=2.4cm]{geometry}
\usepackage{amsmath,amssymb,amsthm}
\usepackage{booktabs}
\usepackage[colorlinks=true,linkcolor=blue,citecolor=blue,urlcolor=blue,
pdftitle={A Decay-Safe Gauged U(1) Origin of the Z17 Spin(10) Axion Model:
Minimal Complete-Pair Anomalons and Two-Loop Audit v20},
pdfauthor={Joel Ayala-Baez}]{hyperref}
\newtheorem{theorem}{Theorem}
\newtheorem{proposition}{Proposition}
\usepackage{placeins}

\title{\bf A Decay-Safe Gauged $U(1)$ Origin of the $\mathbb{Z}_{17}$
Spin(10) Axion Model\\ \large Minimal Complete-Pair Anomalons and Mass-Dependent
Two-Loop Audit}
\author{Joel Ayala-Baez\\
\small Independent Researcher --- \texttt{jayalabaez@gmail.com}}
\date{2 August 2026}

\newcommand{\ueV}{\,\mu\mathrm{eV}}
\newcommand{\GeV}{\,\mathrm{GeV}}
\newcommand{\fa}{f_a}
\newcommand{\ma}{m_a}
\newcommand{\Zs}{\mathbb{Z}_{17}}

\begin{document}
\maketitle

\begin{abstract}
\sloppy
In the SO(10)$\times$U(1)$_{\rm PQ}$ axion model of Ernst, Ringwald and
Tamarit, gauging $\Zs=U(1)_{\rm PQ}\bmod17$ requires anomaly-cancelling
matter. For arbitrary vector-like $16\oplus\overline{16}$ spectator
pairs massed by the PQ singlet $S$, the mixed anomaly alone forces
$k\equiv5\pmod{17}$, independent of their individual charges. Within
the identical-pair ansatz all discrete anomaly classes then select the
unique charge residues $\{2,11\}\pmod{17}$. The minimal completion has
five pairs, shifts the
gluonic coefficient to $2N=-68$, and gives
$N_{\rm DW}^{\rm cover}=17$, $\fa=v_S/17$, unit physical domain-wall
number and $E/N=8/3$. The benchmark $v_S=M_I=6.31\times10^{11}\GeV$,
$y_s=\sqrt2$ predicts $\ma=153.5\pm2.0\ueV$,
$\nu=37.11\pm0.49$~GHz and
$g_{a\gamma\gamma}=(2.335\pm0.125)\times10^{-14}\GeV^{-1}$.

At the discrete-EFT level we distinguish a local
PQ-breaking operator from a spurion set that can close into a vacuum
amplitude. Although $\Zs$ permits fermionic operators of dimensions 6, 9,
10 and 12, the renormalizable theory preserves an
exact spectator vector number $V$. An intentionally over-complete
enumeration of all monomials satisfying Lorentz parity, the Spin(10)
centre and $\Zs$ excludes a PQ-sensitive, $V=0$ closure through
$P=\sum_i(d_i-4)=11$. The bound is saturated at $P=12$ by four
dimension-6 portals plus one PQ-conserving dimension-8 vector-breaking
spurion. An explicit 32-dimensional Clifford-algebra and Grassmann audit
finds nonzero Spin(10) and Lorentz contractions, verifies Fermi statistics,
and identifies a connected two-loop graph. Wilsonian NDA gives
$|\Delta\bar\theta|\lesssim9.04\times10^{-28}|C_{\rm eff}|$. The extra
$\mathbb Z_3$ proposed in v15 is unnecessary and is removed.

We then construct a decay-safe anomaly-free origin of the discrete
symmetry. Gauge $U(1)_X$, break it with a PQ-neutral scalar $\Phi_{17}$,
and add three complete pairs with charges
$(1,16)$, $(14,3)$ and $(1,-18)$. They contribute
$(+34,+272,+16592)$ to the mixed Spin(10), gravitational and cubic
anomalies, exactly cancelling the light sector, while their accidental-PQ
anomaly vanishes. Each pair has a renormalizable portal to ordinary
$16_F$ matter through an already-present $10_H$ or $S$. Mixed-anomaly
cancellation forces the number of single-$\Phi$-massed pairs to be odd;
one pair would require roots of $t^2-17t+76$, whose discriminant is
$-15$. Thus three is minimal in this complete-pair-only ansatz (with no
Spin(10)-singlet anomalons), and exact
enumeration finds the displayed triple uniquely in the stated
existing-field portal basis.

The decay portals change, rather than invalidate, the quality audit. A
heavy-inclusive necessary-condition enumeration excludes a vacuum closure
through $P=7$ and is saturated at $P=8$ by four dimension-five portals
plus one dimension-eight spectator spurion. An explicit nonzero connected
two-loop graph produces $\Phi^4(S^\dagger)^{18}(10_H^\dagger)^2$.
Its repeated-pole integral retains all three masses and gives
$2|A_8|/\chi=6.04\times10^{-47}|C_8|$ at
$v_\Phi=10^{17}\GeV$. The direct dimension-21 scalar operator remains
the largest explicitly computed unit-coefficient term,
$2|A|/\chi=4.60\times10^{-43}$. A normalized decay portal as small as
$3.1\times10^{-20}$ gives a lifetime below one second; every 16 component
has a nonzero $10_H$ channel in the explicit Clifford representation.
The construction is anomaly-free as an effective field theory, not a string
completion. Perturbativity to the reduced Planck scale is not established:
the present one-loop Spin(10) envelope reaches a Landau pole below that scale.
Flavour/Clebsch fits, hierarchy stabilization, the full scalar vacuum and the
string network remain open. Three fail-closed
engines perform 65, 59 and 42 checks; the release adds independent charge,
decay and loop-quadrature tests.
\end{abstract}

\section{Model and conventions}
Fields and PQ charges: $16_F$ ($+1$, three generations), $10_H$ ($-2$),
$\overline{126}_H$ ($-2$), $210_H$ (0), $S$ ($+4$) \cite{ERT}; the axion
is intended to be dominantly $\arg S$. The previously proposed 54-channel
$\overline{126}_H^2 10_H^2 S^2$ locking term vanishes on the selected
$\Delta_R^2$ projection, so the complete phase Hessian remains open. At low
energy we impose $\Zs := \mathrm{U(1)}_{\rm PQ}\bmod 17$; Section~\ref{sec:uv}
realizes it as the remnant of a primitive gauged $U(1)_X$. Spectators: $k$
identical $16_s\oplus\overline{16}_s$ pairs
with $\Zs$ charges $(a,b)$, masses from $S\,16_s\overline{16}_s$, i.e.\
thresholds $M_s = y_s v_S/\sqrt2$. Discrete-gauge protection of axions
has a long history \cite{BGW,DPT,DiLuzio,BDM}; the specific result
below concerns this particular $\Zs$ extension of the ERT model.

\section{The theorem}
\begin{theorem}[minimality and residue uniqueness, identical-pair ansatz]
The $\Zs$ anomaly classes (mixed SO(10)$^2$, gravitational, cubic
\cite{IR,Hsieh}) all vanish modulo 17 if and only if
$k \equiv 5 \pmod{17}$ and $\{a,b\} \equiv \{2,11\} \pmod{17}$. The
minimal solution therefore contains five pairs, and the charge
\emph{residues} are unique up to interchange; $(2,-6)$ is one integer
lift.
\end{theorem}

\begin{proof}
The mass operator forces $a + b \equiv -4 \equiv 13$. The mixed class
is $6 + 2k(a{+}b) \equiv 6 - 8k$; with $8^{-1} \equiv 15$, vanishing
requires $k \equiv 90 \equiv 5$, i.e.\ $k = 5 + 17m$. The gravitational
class $48 + 16k(a{+}b) \equiv 14 + 4k$ vanishes automatically on the
same set. The cubic class uses
$a^3 + b^3 = (a{+}b)^3 - 3ab(a{+}b) \equiv 4 - 5ab$, giving
$48 + 80(a^3{+}b^3) \equiv 11 - 9ab \equiv 0$, i.e.\ $ab \equiv 5$. The
residue pair is then the root set of $x^2 - 13x + 5 \equiv 0
\pmod{17}$, whose only roots are $2$ and $11$; an exhaustive residue
scan confirms uniqueness.
\end{proof}

The five-pair lower bound is stronger than the theorem's identical-pair
wording. For arbitrary residues $(a_i,13-a_i)$, the mixed class is
\begin{equation}
 6+2\sum_{i=1}^{k}\bigl[a_i+(13-a_i)\bigr]
 \equiv 6+9k\pmod{17},
\end{equation}
so every charge assignment requires $k\equiv5\pmod{17}$. Direct
enumeration gives no solution for $k\leq4$ and $83\,232$ ordered cubic
solutions among the $17^5$ tuples at $k=5$. Thus pair-count minimality is
charge-independent, whereas the residue uniqueness $\{2,11\}$ is
specifically an identical-pair result. The general cubic class used here is
$48+16\sum_i[a_i^3+(13-a_i)^3]$; the alternate engine incorrectly used
80 per pair (valid only after five identical pairs), although its solution
count happened to remain unchanged.

The five pairs shift the gluonic coefficient to
$2N = 2[6 + 10(a{+}b)] = -68$, so
\begin{equation}
N_{\rm DW}^{\rm cover} = \frac{|2N|}{X_S} = \frac{68}{4} = 17,
\qquad
\fa \longrightarrow \frac{v_S}{17},
\end{equation}
in the discrete EFT, or equivalently the $v_\Phi\gg v_S$ limit of
Section~\ref{sec:uv}; the exact continuous-theory projection is given in
Eq.~\eqref{eq:axionprojection}. The ratio $E/N = 8/3$ is preserved
(complete SO(10) multiplets). All seventeen
covering-space minima form one gauge orbit
($4\times13\equiv1\bmod17$): the physical domain-wall number is one,
and the one-wall object is the string sector $(\ell,n) = (13,-3)$ with
$S$-winding $2\pi/17$ (we reserve $k$ for the pair multiplicity).

\section{Benchmark and inputs}
\label{sec:bench}
One-loop running of $\alpha_1^{-1}(M_Z)=59.02$ (GUT-normalised),
$\alpha_2^{-1}(M_Z)=29.57$, $\alpha_s(M_Z)=0.1179$ with 2HDM
coefficients $(21/5,-3,-7)$ below $M_I$, Pati--Salam coefficients
$(-7/3,2,26/3)$ above it, and matching
$\alpha_{2R}^{-1} = \tfrac53\alpha_1^{-1} - \tfrac23\alpha_4^{-1}$
gives $M_I = 6.31\times10^{11}\GeV$,
$M_{\rm GUT} = 9.92\times10^{15}\GeV$,
$\alpha_{\rm GUT}^{-1} = 37.31$. \emph{Gauge unification does not fix
the singlet scale \cite{ERT}; we separately impose the benchmark
identifications} $v_S = M_I$ \emph{and} $y_s = \sqrt2$ (so
$M_s = M_I$). Then, with $\chi^{1/4} = 75.5(5)$~MeV and the QCD photon
coefficient $1.92(4)$ from Grilli di Cortona et al.\ \cite{GrillidiCortona},
\begin{align}
\ma &= \frac{\chi^{1/2}}{\fa} = 153.5 \pm 2.0\ueV,
\qquad
\nu = \frac{\ma}{h} = (37.11\pm0.49)\ \mathrm{GHz},\\
g_{a\gamma\gamma} &= \frac{\alpha}{2\pi\fa}
\left|\frac{E}{N}-1.92\right|
= (2.335\pm0.125)\times10^{-14}\GeV^{-1},
\end{align}
where the mass and frequency uncertainties propagate the susceptibility
error and the coupling uncertainty propagates the $1.92(4)$ coefficient.
The compatible determination $\chi^{1/4}=75.6(1.8)(0.9)$~MeV
\cite{Borsanyi} would widen the mass uncertainty to about $8\ueV$.
The spectator shift is
$\Delta\alpha_{\rm GUT}^{-1} = -\tfrac{40/3}{2\pi}
\ln(M_{\rm GUT}/M_s) = -20.5$ at the benchmark $y_s$; other $y_s$
values shift it (e.g.\ $-26.1$ at $y_s = 0.1$), so $y_s$ is an explicit
input of the proton estimate.

\section{Axion quality: local operators and vacuum closure}
\label{sec:quality}
At the reduced Planck scale we assume every local operator allowed by
the gauge group is present, with a dimensionless Wilson coefficient of
order unity, and that no intermediate threshold carries additional
explicit PQ breaking. This assumption is essential: heavy fields can
modify naive Planck power counting \cite{Bonnefoy}.

\subsection*{Scalar sector and the local-operator trap}
The $10_H$ and $\overline{126}_H$ tensors have an odd number of SO(10)
vector indices, whereas $210_H$ has an even number. Deltas and the
rank-ten epsilon tensor contract an even total number of vector indices,
so an invariant contains an even number of odd-index tensors. Including
$S$ (charge $+4$), every scalar invariant therefore has
$Q_{\rm PQ}=0\pmod4$. Combining this with $\Zs$ gives
$|Q_{\rm PQ}|\geq\mathop{\rm lcm}(4,17)=68$, saturated by
$S^{17}/M_{\rm Pl}^{13}$. At the benchmark, the one-sided coefficient is
$(v_S/\sqrt2)^{17}/(M_{\rm Pl}^{13}\chi)=3.24\times10^{-37}$.
With $V=A e^{ia/f_a}+{\rm h.c.}$ the worst-phase shift is
\begin{equation}
 \frac{|\Delta\bar\theta|_{S^{17}}}{|c_{17}|}
 = \frac{2(v_S/\sqrt2)^{17}}{M_{\rm Pl}^{13}\chi}
 =6.47\times10^{-37},
 \qquad M_{\rm Pl}=2.435\times10^{18}\GeV .
\label{eq:scalarquality}
\end{equation}

Fermions invalidate any proof that stops there; multi-fermion operators
are known to matter in axion-quality analyses \cite{BDMquality}. Write
$F=16_F$, $s=16_s$ and $b=\overline{16}_s$, with
$Q_{\rm PQ}=(1,2,-6)$. SO(10)$\times\Zs$ permits, among others,
\begin{align}
{\cal O}_6^-&=\frac{(S^\dagger)^3(Fb)_1}{M_{\rm Pl}^{2}},
 &(d,Q,V)&=(6,-17,-1),\nonumber\\
{\cal O}_8^0&=\frac{(S^\dagger)^2[(ss)_{10}]^2}{M_{\rm Pl}^{4}},
 &(d,Q,V)&=(8,0,+4),\nonumber\\
{\cal O}_9&=\frac{(S^\dagger)^5(bb)_{10}10_H}{M_{\rm Pl}^{5}},
 &(d,Q,V)&=(9,-34,-2),\nonumber\\
{\cal O}_{10}&=\frac{10_H^\dagger[(bb)_{10}]^3}{M_{\rm Pl}^{6}},
 &(d,Q,V)&=(10,-34,-6),\nonumber\\
{\cal O}_{12}&=\frac{(S^\dagger)^6(Fs)_{10}(ss)_{10}}
 {M_{\rm Pl}^{8}},
 &(d,Q,V)&=(12,-17,+3),
\label{eq:operators}
\end{align}
where $V(s)=+1$, $V(b)=-1$ and all ERT fields have $V=0$.
The representation products used in these explicit contractions are
standard, in particular $16\otimes16=10_s\oplus120_a\oplus126_s$
\cite{Djouadi}. The dimension-10 and -12 terms are regression cases
missed by a bilinear-only enumeration.

Nor does ${\cal O}_9$ by itself generate the previously claimed term
linear in its Majorana insertion. For one Dirac pair,
\begin{equation}
 {\cal M}=\begin{pmatrix}0&M_s\\M_s&\delta e^{i\phi}\end{pmatrix},
 \quad {\rm tr}({\cal M}^\dagger{\cal M})=2M_s^2+|\delta|^2,
 \quad \det({\cal M}^\dagger{\cal M})=M_s^4 .
\label{eq:massmatrix}
\end{equation}
Both singular values, hence the one-loop Coleman--Weinberg potential,
are independent of $\phi$. A local PQ-breaking operator is therefore
not automatically a one-insertion axion potential.

\subsection*{Finite multi-spurion proof}
The renormalizable Lagrangian preserves the spectator vector number
$U(1)_V$: the mass $Ssb$ is neutral and there is no dimension-four or
lower $Q_{\rm PQ}=0$, $V\neq0$ interaction. Consequently a vacuum graph
must have total $V=0$. Each Planck insertion has $Q_{\rm PQ}=17m$; after
its fermions are closed with PQ-conserving vertices, scalar index parity
also requires total $Q_{\rm PQ}=0\pmod4$. A PQ-sensitive target thus has
\begin{equation}
 \sum_i V_i=0,\qquad \sum_i m_i\neq0,\qquad
 \sum_i m_i=0\pmod4 .
\label{eq:closureconditions}
\end{equation}

The supplied engine enumerates a deliberate \emph{superset} of local
monomials: $F,s,b$ and all conjugates; $S,S^\dagger$; charge-$\mp2$
tensor scalars; even Weyl count; trivial Spin(10)-centre charge; and
$Q_{\rm PQ}=0\pmod{17}$. It allows mixed-chirality contractions that may
not exist, so failure to find a closure is a conservative exclusion.
Neutral $210_H$, derivatives and gauge strengths can be omitted because
they only raise dimension. For
\begin{equation}
 P:=\sum_i(d_i-4)\leq11,
\end{equation}
every insertion has $d_i\leq15$; enumeration through dimension 15 is
therefore finite and complete for this exclusion. No state satisfying
Eq.~\eqref{eq:closureconditions} occurs through $P=11$. At $P=12$ the
bound is saturated by the actual Spin(10) set
\begin{equation}
 4{\cal O}_6^-+{\cal O}_8^0:\qquad
 (Q_{\rm PQ},V,P)=(-68,0,12).
\label{eq:certificate}
\end{equation}
The necessary-condition search was independently reimplemented directly
from weak compositions of field multiplicities; it reproduces all twenty
low-cost $(P,Q_{\rm PQ}/17,V)$ frontier classes and the same absence result
through $P=11$ without importing the main enumeration module.

\subsection*{Explicit Spin(10), Lorentz and graph audit}
Because spinors are representations of the double cover, this audit uses
Spin(10) explicitly (``SO(10)'' elsewhere is the standard shorthand).
Let $\Gamma^a$ be a $32\times32$ Euclidean Clifford representation and
$C=\Gamma^2\Gamma^4\Gamma^6\Gamma^8\Gamma^{10}$. On either chiral
16-dimensional subspace, the ten matrices
\begin{equation}
 T^a_\pm=(C\Gamma^a)|_{16_\pm},\qquad
 (T^a_\pm)^T=T^a_\pm,\qquad
 \sum_{ij}(T^a_+)^*_{ij}(T^b_+)_{ij}=16\delta^{ab}
\label{eq:cliffordcertificate}
\end{equation}
are symmetric, so the same-chirality Weyl bilinear
$\epsilon_{\alpha\beta}X_i^\alpha(T^a_\pm)_{ij}Y_j^\beta$
has the required Fermi statistics for the $10_s$ channel. Sparse
Grassmann evaluation gives nonzero polynomials for every operator in
Eq.~\eqref{eq:operators}; in particular ${\cal O}_8^0$ has 240 nonzero
canonical monomials even when its four $s$ fields share a species.

For the central four-spinor tensor
$K_{ijkl}=\sum_a(T^a_+)_{ij}(T^a_+)_{kl}$, the explicit representation
gives 640 nonzero entries and
\begin{equation}
 \sum_{ijkl}|K_{ijkl}|^2=2560,
 \qquad K^*_{ijkl}(T^1_+)_{ij}(T^1_+)_{kl}=256,
 \qquad L_{\rm Lorentz}=4 .
\label{eq:graphcertificate}
\end{equation}
Thus the certificate is not merely a centre-charge candidate: its actual
group and Lorentz contractions are nonzero. Four $s$--$b$ same-chirality
propagators (the four $s$ legs may be assigned four of the five spectator
copies) and two ordinary-family $10_H$-channel chirality flips close the
twelve fermions. With five compact vertices and six internal fermion lines,
$L=6-5+1=2$. The resulting scalar operator is proportional to
\begin{equation}
 (S^\dagger)^{18}(10_H^\dagger\!\cdot10_H^\dagger),
 \qquad Q_{\rm PQ}=-68 .
\label{eq:scalarphase}
\end{equation}
A special Higgs alignment could make the dot product vanish, which would
only strengthen the quality bound by moving the first nonzero vacuum term
to higher order.

Treating all masses and scalar insertions as no larger than the matching
scale $M_s$ and dropping loop suppressions gives the Wilsonian-NDA upper bound
\begin{equation}
 \frac{|\Delta\bar\theta|}{|C_{\rm eff}|}
 \lesssim \frac{2M_s^4}{\chi}
 \left(\frac{M_s}{M_{\rm Pl}}\right)^{12}
 =9.04\times10^{-28} .
\label{eq:qualitybound}
\end{equation}
Here $C_{\rm eff}$ absorbs convention-dependent tensor normalisations,
Yukawas and flavour factors. For the explicit graph, retaining the two
loop factors and using $|v_{10}|\leq246\GeV$ yields the diagnostic
\begin{equation}
 \frac{|\Delta\bar\theta|_{P=12}}{|C_{\rm eff}|}
 \lesssim 9.04\times10^{-28}(16\pi^2)^{-2}
 \left(\frac{246\GeV}{M_s}\right)^2
 =5.50\times10^{-51},
\label{eq:refinedgraph}
\end{equation}
before any additional flavour suppression. Equation~\eqref{eq:qualitybound},
which deliberately omits these small factors, remains the advertised
conservative result.

\begin{proposition}[quality under the stated EFT assumptions]
At the benchmark $M_s=v_S=6.31\times10^{11}\GeV$, the minimal
SO(10)$\times\Zs$ completion satisfies the neutron-EDM requirement
$|\Delta\bar\theta|<10^{-10}$ provided
$|C_{\rm eff}|<1.1\times10^{17}$. In particular, order-one Wilson
coefficients are safe; the leading conservative bound is
Eq.~\eqref{eq:qualitybound}, not the scalar-only value in
Eq.~\eqref{eq:scalarquality}.
\end{proposition}

\section{An anomaly-free continuous origin of
\texorpdfstring{$\mathbb Z_{17}$}{Z17}}
\label{sec:uv}
A discrete gauge symmetry can arise as the remnant of a broken continuous
gauge group \cite{KW,BanksDine,Ibanez93}. Introduce a complex Spin(10)
singlet $\Phi$ with $(X,{\rm PQ})=(17,0)$ and gauge a primitive
$U(1)_X$. The light-field $X$ charges are the integer PQ lifts used above:
$X(F,s,b,S,10_H,\overline{126}_H,210_H)=(1,2,-6,4,-2,-2,0)$.
Replace the relic-unsafe v19 anomalons by the left-Weyl pairs
\begin{equation}
 (P,\overline P)\sim(16_{1},\overline{16}_{16}),\qquad
 (Q,\overline Q)\sim(16_{14},\overline{16}_{3}),\qquad
 (R,\overline R)\sim(16_{1},\overline{16}_{-18}).
\label{eq:uvfields}
\end{equation}
Their accidental-PQ charges are respectively $(1,-1)$, $(-3,3)$ and
$(1,-1)$. Renormalizable masses and decay portals are
\begin{equation}
 \Phi^\dagger(y_P P\overline P+y_QQ\overline Q)
 +y_R\Phi R\overline R
 +\lambda_P P F10_H+\lambda_R R F10_H
 +\lambda_Q\overline QF S^\dagger+{\rm h.c.}
\label{eq:uvmasses}
\end{equation}
Every $X$ and PQ charge sum in Eq.~\eqref{eq:uvmasses} vanishes exactly.
The two added $16_1$ fields share the ordinary-family gauge charges. A
generic rank-two $\Phi$ mass matrix pairs two linear combinations of the
five $16_1$ fields with $\overline{16}_{16}$ and
$\overline{16}_{-18}$, leaving exactly three chiral light combinations.
Thus the labels $F,P,R$ denote a convenient UV basis, not a claim that
the heavy--light mass matrix is diagonal.

With $T(16)=2$ and $\dim16=16$, the continuous anomaly audit is
\begin{center}
\begin{tabular}{lrrr}
\toprule
sector & $A_{10^2X}$ & $A_{{\rm grav}^2X}$ & $A_{X^3}$\\
\midrule
light fields & $-34$ & $-272$ & $-16592$\\
three pairs in Eq.~\eqref{eq:uvfields} & $+34$ & $+272$ & $+16592$\\
\midrule
total & $0$ & $0$ & $0$\\
\bottomrule
\end{tabular}
\end{center}
The cubic entries before the common factor 16 are
$4097+2771-5831=1037$. The heavy accidental-PQ mixed anomaly is also zero,
so the gauge-anomaly coefficients are unchanged. This does not by itself
fix the regular light-fermion current. For a complete singlet-VEV mass block
with submatrices $(A,B,C,D)$, where $C$ includes the allowed $SQ\overline R$
portal, let $B_\ell(\alpha)$ embed the three canonically normalized light
modes. Then
\begin{equation}
 Q_{\rm proj}=B_\ell^\dagger Q_{\rm UV}B_\ell
 =\mathbf1-4W,\qquad
 {\cal A}_\alpha=iB_\ell^\dagger\frac{dB_\ell}{d\alpha}=4W.
\label{eq:portalcurrent}
\end{equation}
The sum $Q_{\rm proj}+{\cal A}_\alpha=\mathbf1$ is a moving-coordinate
identity and is basis dependent; it does not erase the physical projected
current. General portal matrices make $W$ non-universal, and misalignment
with the SM Yukawa matrices can generate flavour-changing axion currents.
A full matching therefore requires the portal tensors and component Yukawa
matrices, even though the total anomaly is unchanged.

There is a short minimality proof for an anomalon sector built only from
complete pairs, with no Spin(10)-singlet anomalons. If every
$16\oplus\overline{16}$ pair receives a renormalizable mass from one
$\Phi$ or $\Phi^\dagger$, its two $X$ charges sum to $\pm17$. Cancellation
of the mixed anomaly requires $n_+-n_-=1$, hence an odd number of pairs.
For one pair, $x+y=17$ and cubic cancellation require
$x^3+y^3=1037$, or $xy=76$. The roots would satisfy
$t^2-17t+76=0$, whose discriminant is $-15$. One pair is impossible,
whereas Eq.~\eqref{eq:uvfields} supplies three; three is therefore minimal
under these assumptions. Exact enumeration of portals constructed from
the already-present $10_H$, $S$ and $210_H$ finds the displayed triple as
the unique three-pair solution. This is not a classification involving
arbitrary Spin(10) representations, singlet anomalons or multi-scalar mass
operators.

An over-catalogue of every dimension-four-or-lower monomial satisfying
Lorentz parity, the Spin(10) centre and exact $U(1)_X$ contains no term
that breaks accidental PQ or the light-spectator vector number. The heavy
pair numbers are explicitly broken by Eq.~\eqref{eq:uvmasses}, as required
for decay. Hence PQ remains accidental but no exact anomalon number makes a
heavy state stable. PQ is independent of the gauged generator because
$\Phi$ has $X=17$ but PQ zero.
Writing the canonically normalized phases as $(a_\Phi,a_S)$, define
\begin{equation}
 D=\sqrt{(17v_\Phi)^2+(4v_S)^2},\qquad
 \widehat a=\frac{-4v_Sa_\Phi+17v_\Phi a_S}{D},\qquad
 f_a=\frac{v_Sv_\Phi}{D}.
\label{eq:axionprojection}
\end{equation}
The first combination is orthogonal to the gauge direction. For
$v_\Phi=10^{17}\GeV$ the correction to $v_S/17$ is
$-1.10\times10^{-12}$.

The exact coefficient converting the $S$ phase into $\widehat a/f_a$ is
\begin{equation}
 \xi=\frac{17v_\Phi^2}{(17v_\Phi)^2+(4v_S)^2}
 =0.0588235294116349
\label{eq:xiexact}
\end{equation}
at the frozen v20 point. If the physical projected current is aligned as
$Q_{\rm proj}=\mathbf1$ in the SM mass bases, standard axial-basis matching
gives the conditional tree benchmark
\begin{equation}
 C^0_{u,c,t}=\xi\cos^2\!\beta,\qquad
 C^0_{d,s,b}=C_e=\xi\sin^2\!\beta.
\label{eq:fermiontreeexact}
\end{equation}
Using central low-energy hadronic coefficients \cite{GrillidiCortona},
\begin{align}
 C_p&=-0.47+\xi(0.8645\cos^2\!\beta-0.437\sin^2\!\beta),\nonumber\\
 C_n&=-0.02+\xi(-0.4055\cos^2\!\beta+0.833\sin^2\!\beta).
\label{eq:nucleoncentral}
\end{align}
The constants in Eq.~\eqref{eq:nucleoncentral} carry correlated hadronic
uncertainties of order $0.03$ and a precision result requires a declared
threshold/RG convention. At the constrained benchmark $\tan\beta=1.5$ the
central values are
\begin{equation}
 (C_e,C_p,C_n)=(0.04072398,-0.47214932,0.00658371).
\end{equation}
These are not unique numerical v20 predictions. They assume current/Yukawa
alignment, while Eq.~\eqref{eq:portalcurrent} shows that general portals can
change the physical current. Moreover, the corrected flavour analysis must
use Takagi diagonalization of the complex-symmetric Majorana matrix and
$U_{\rm PMNS}=U_e^\dagger U_\nu$. After this correction the previous
$\tan\beta=1.5$ and high-$\tan\beta$ minima are invalid; the current
$v_R=v_S$ profile has no point with $\chi^2<30$. Precision matching also
requires threshold/RG evolution and correlated hadronic uncertainties.

The large charges make the Abelian running a real constraint. Including
all fields above $v_\Phi$ gives
\begin{equation}
 b_X=\frac23\sum_{\rm Weyl}d(R)X^2+
 \frac13\sum_{\rm scalar}d(R)X^2
 =\frac23(15840)+\frac13(849)=10843.
\end{equation}
Requiring the one-loop Landau pole to exceed
$M_{\rm Pl}=2.435\times10^{18}\GeV$ at $v_\Phi=10^{17}\GeV$ requires
$g_X(v_\Phi)<0.0478$; $g_X=0.04$ gives
$\Lambda_{\rm LP}=9.47\times10^{18}\GeV$. Complete Spin(10) multiplets
do not change one-loop relative unification below their common threshold.
Using $T(10,126,210)=(1,35,56)$ \cite{Chang} and conservatively treating
even the real $210_H$ as complex gives $b_{10}\leq80/3$ above $v_\Phi$.
An older diagnostic reset $\alpha_{10}(v_\Phi)=1/40$ would remain
perturbative, but that reset is inconsistent with the model's own
spectator-corrected $\alpha_{\rm GUT}^{-1}\simeq16.81$. Continuous
one-loop running from that value is \emph{not} weakly coupled to
$M_{\rm Pl}$ under these betas (the conservative envelope hits a Landau
pole below $M_{\rm Pl}$). A full two-loop threshold analysis is therefore
required before claiming Planck-safe Spin(10) perturbativity.

The decay statement is quantitative. For one open normalized two-body
channel with final mass $m_\phi$,
\begin{equation}
 \Gamma=\frac{|\lambda|^2M}{32\pi}\Bigl(1-\frac{m_\phi^2}{M^2}\Bigr)^2
 \le\frac{|\lambda|^2M}{32\pi},\qquad
 M=\frac{v_\Phi}{\sqrt2}.
\end{equation}
The massless formula is therefore an upper benchmark on the width. At the
benchmark, $\lambda=10^{-8}$ and $m_\phi\to0$ gives
$\tau=9.36\times10^{-24}$~s, while $\tau<1$~s in that massless limit
requires only $|\lambda|>3.06\times10^{-20}$. In the explicit Clifford basis,
$\sum_{a,j}(C\Gamma^a)_{ij}(C\Gamma^a)^*_{kj}=10\delta_{ik}$, proving
that every component of $P$ and $R$ has a nonzero $10_H$ channel. The
$\overline QFS^\dagger$ singlet contraction is component-diagonal. This
removes the exact-stable-relic obstruction; a broken-phase Clebsch and
flavour fit is still required for precision widths.

Finally, the portal
$|\Phi|^2|S|^2$ must be small enough to preserve $v_\Phi\gg v_S$.
These are explicit perturbativity and hierarchy assumptions, so this is a
field-theory EFT valid only below its first strong-coupling or UV-completion
threshold. The present running does not establish a weakly coupled completion
to the Planck cutoff, let alone an ultimate quantum-gravity construction.

\section{Heavy-threshold closure and full loop kernels}
\label{sec:matching}
Heavy fields can invalidate naive axion-quality power counting
\cite{Bonnefoy}, so the operator search was repeated with all fields in
Eq.~\eqref{eq:uvfields} and their conjugates. Because the decay vertices
break the heavy pair numbers, only the exact light-spectator number $V_s$
is imposed. The resulting minimum-dimension frontier has no PQ-sensitive,
$V_s=0$ closure through $P=7$. At $P=8$ the necessary-condition bound is
saturated by actual Spin(10) invariants:
\begin{align}
 \widetilde{\cal O}_5&=\frac{(S^\dagger)^2(Qb)_1}{M_{\rm Pl}},
 &(Q_{\rm PQ},V_s,P)&=(-17,-1,1),\nonumber\\
 \widetilde{\cal O}_8&=\frac{(S^\dagger)^2
 [(ss)_{10}]^2}{M_{\rm Pl}^4},
 &(Q_{\rm PQ},V_s,P)&=(0,+4,4).
\label{eq:p8ops}
\end{align}
The combination $4\widetilde{\cal O}_5+\widetilde{\cal O}_8$ closes with
four chirality-flipping $s$--$b$ propagators generated by $Ssb$, four
$Q$--$\overline Q$ propagators generated by
$\Phi^\dagger Q\overline Q$, four $\overline QFS^\dagger$ decay vertices
and two conjugate ordinary $FF10_H$ Yukawas. The left--left propagators
carry the conjugate mass spurions, $S^\dagger$ and $\Phi$, respectively.
Counting the five Planck vertices, four decay vertices and two ordinary
Yukawas gives a connected graph with 12 internal fermion lines, 11
vertices and two loops. It produces
\begin{equation}
 \Phi^4(S^\dagger)^{18}(10_H^\dagger)^2,
 \qquad (X,Q_{\rm PQ},P)=(0,-68,8).
\end{equation}
The explicit Clifford representation gives
$(T^a)^*_{ij}(T^b)_{ij}=16\delta^{ab}$ and the inherited four-bilinear
two-component contraction, including both kinetic numerators, equals $-2$
per factorized loop (hence 4 for their product), so this saturation graph
is nonzero.

The original v17 two-loop graph also survives, but its four
${\cal O}_6^-$ portals each carry $X=-17$. Exact gauge invariance replaces
each by $\Phi{\cal O}_6^-/M_{\rm Pl}$; the graph is therefore $P=16$.
The full zero-momentum integral can be written without dimensional
estimates. For degenerate spectator mass $M_s$ and a family chirality
mass $m$,
\begin{align}
 I_3(M_s,m)&=\int\frac{d^4p_E}{(2\pi)^4}
 \frac{1}{(p^2+M_s^2)^2(p^2+m^2)}\nonumber\\
 &=\frac{1}{16\pi^2M_s^2}
 \frac{1-r+r\ln r}{(1-r)^2},\qquad r=\frac{m^2}{M_s^2},
 \label{eq:triangle}\\
 J(M_s,m)&=M_s^2m\,I_3(M_s,m),\nonumber\\
 |A_{2\ell}^{\rm EFT}|&=|C_{2\ell}|\frac{s_0^{14}}{M_{\rm Pl}^{12}}
 J(M_s,m_1)J(M_s,m_2),\qquad s_0=\frac{v_S}{\sqrt2}.
\label{eq:full2loop}
\end{align}
For unequal masses the code evaluates the symmetric divided difference
\begin{equation}
 I_3=\frac{1}{16\pi^2}\sum_x
 \frac{x\ln x}{\prod_{y\ne x}(x-y)}.
\end{equation}
Equation
\eqref{eq:full2loop} is finite and factorizes because the zero-momentum
four-fermion vertex leaves two independent loop momenta. The normalized
coefficient $C_{2\ell}$ contains the unknown Planck Wilson tensors,
flavour contractions, the nonzero group/Lorentz normalization and RG
matching; the EFT cannot predict it.

For the $P=8$ graph, each of the two factorized loops has the independently
evaluated kernel
\begin{equation}
 K(M_H,M_s,m_f)=M_H^2M_s^2\int\frac{d^4p_E}{(2\pi)^4}
 \frac{p^2}{(p^2+M_H^2)^2(p^2+M_s^2)^2(p^2+m_f^2)^2}.
\label{eq:chain}
\end{equation}
High-precision repeated-pole partial fractions retain all three masses and
are independently checked by log-coordinate quadrature. The one-sided
amplitude is
\begin{equation}
 |A_8|=|C_8|\frac{s_0^{14}h^2}{M_{\rm Pl}^{8}}
 K(M_H,M_s,m_f)^2.
\label{eq:fullp8}
\end{equation}
Here $C_8$ contains the four decay Yukawas, two family Yukawas and the
normalized Wilson, group, flavour and RG contraction. With $M_s=v_S$,
$M_H=v_\Phi/\sqrt2$, $m_f=h=246\GeV$ and $v_\Phi=10^{17}\GeV$,
$16\pi^2M_H^2K=0.9999999965$. The results per unit normalized
coefficient are
\begin{center}
\begin{tabular}{lcc}
\toprule
term & $|A|/\chi$ & worst phase $2|A|/\chi$\\
\midrule
v17 two-loop EFT graph, before $\Phi$ dressings & $2.15\times10^{-53}$ &
$4.30\times10^{-53}$\\
same graph in $U(1)_X$ ($P=16$) & $1.53\times10^{-59}$ &
$3.06\times10^{-59}$\\
decay-threshold graph ($P=8$) & $3.02\times10^{-47}$ &
$6.04\times10^{-47}$\\
direct $\Phi^4(S^\dagger)^{17}/M_{\rm Pl}^{17}$ &
$2.30\times10^{-43}$ & $4.60\times10^{-43}$\\
\bottomrule
\end{tabular}
\end{center}
The direct dimension-21 scalar operator dominates these explicitly
computed terms and still lies $2.17\times10^{32}$ below the
$10^{-10}$ criterion for a unit coefficient. This numerical margin is
conditional on order-one normalized Wilson contractions; it is not a
measurement or a derivation of quantum-gravity coefficients.

An extra vector-like spectator $\mathbb Z_3$ is neither needed nor used.
Indeed ${\cal O}_{10}$ and ${\cal O}_{12}$ survive it, falsifying the
earlier claim that all fermionic breakers move to dimension 19. Such an
exact triality would also make the lightest triality-charged state stable
unless further structure were added. The rank-six rephasing diagnostic
in Appendix~A remains correct but is not the quality proof. The continuous
origin is supplied in Section~\ref{sec:uv}; its cosmic-string spectrum
remains open.

\subsection*{Adversarial merge audit of the alternate anomalon lift}
The alternate charges
$16_{1}\oplus\overline{16}_{16}$ and singlets
$(33,-16,31,-48)$ do cancel the three continuous anomaly sums exactly.
They do not, however, establish the advertised quality theorem. A physical
axion first requires a global accidental PQ independent of the gauged
generator; one possible repair assigns
$Q_{\rm PQ}(16_1,\overline{16}_{16},\Phi)=(1,-1,0)$.
With that assignment the following actual 10-channel invariants are
allowed:
\begin{align}
 {cal O}_{6}^{A}&=\frac{(S^\dagger)^2
   (\overline{16}_{16}\,\overline{16}_{s,-6})_{10}10_H}{M_{\rm Pl}^2},
 &(Q_{\rm PQ},V_s,P)&=(-17,-1,2),\nonumber\\
 {cal O}_{8}^{s}&=\frac{(S^\dagger)^2
   [(16_{s,2}16_{s,2})_{10}]^2}{M_{\rm Pl}^4},
 &(Q_{\rm PQ},V_s,P)&=(0,+4,4).
\end{align}
Hence $4{\cal O}_{6}^{A}+{\cal O}_{8}^{s}$ is a
$(Q_{\rm PQ},V_s,P)=(-68,0,12)$ closure. The proposed
$16_1 16_F10_H$ vertex breaks anomalon number, and four anomalon-mass
propagators, four such vertices and two ordinary $16_F16_F10_H$ vertices
provide a connected two-loop topology for the remaining fermions. An
end-to-end Grassmann contraction of that alternate topology is still
required. The
scalar-only dimension-21 scan therefore misses the relevant fermionic
threshold; this does not prove that the alternate model is unsafe, but it
does mean its published quality bound and loop amplitude are incomplete.

The same audit finds only the two $\Phi$ mass Yukawas for the four heavy
singlets, not a renormalizable portal to light matter. It also gives
$b_X=8263$ and $g_X(M_{\rm GUT})<0.0417$ for a Landau pole above
$M_{\rm Pl}$. Finally, applying the exact factorized kernel to the
\emph{old} graph at the alternate inputs gives
$|A|/\chi=1.85\times10^{-64}$ after four dressings, rather than
$2.37\times10^{-62}$; the ratio is the earlier 128-fold vev-normalization
envelope. These are the reasons the alternate anomalon sector is not
merged, while its charge-general minimality proof is retained.

\section{Status of cosmology, experiment, and proton decay}

\emph{Cosmology (conditional).} The v20 anomalons carry no exact heavy
number and decay through Eq.~\eqref{eq:uvmasses}; unlike v19, avoiding a
stable heavy relic does not require a low-reheat postulate. This statement
does not replace a Boltzmann calculation if the $U(1)_X$ sector is
thermally restored. Pre-inflationary: $\theta_i = 2.91$ on
the anharmonic tail \cite{KKT}, $H_I \lesssim 9\times10^{5}\GeV$
(UV-conditional \cite{GR25}), corresponding to
$r<1.34\times10^{-17}$ for $A_s=2.1\times10^{-9}$. A primordial
B-mode detection above that ceiling rejects this pre-inflationary
all-dark-matter branch, not the post-inflationary branch. The latter is
conditional on the
$(\ell,n)=(13,-3)$ string's existence and production \cite{LRS}. If a
standard network applied, current simulations disagree on the full-DM
mass: Benabou et al.\ find $45$--$65\ueV$ before wall collapse, with
wall effects potentially extending to $\sim300\ueV$ \cite{Benabou};
Saikawa et al.\ report $95$--$450\ueV$ \cite{Saikawa}. The benchmark
lies inside the latter and inside the extended range of the former;
the nonstandard gauged network makes all such intervals indicative.

\emph{Experiment.} No dark-matter haloscope has scanned $153.5\ueV$: the
recent DALI results sit at $28.5\ueV$ \cite{DALI} (their reported
sensitivity is disputed \cite{DALIcomment}), and no existing experiment
reaches the predicted coupling there; CAST constrains the mass at
$5.8\times10^{-11}\GeV^{-1}$ \cite{CAST}, a factor 2\,479 above.
Direct tests belong to MADMAX \cite{MADMAX} and plasma-haloscope
\cite{ALPHA} programmes; the ORGAN design range also overlaps the target
\cite{ORGAN}. Their design ranges are not completed scans or guaranteed
DFSZ sensitivity. The benchmark location is $36.62$--$37.60$~GHz, while
the Galactic signal width is only about $34$~kHz for
$Q_{\rm halo}=1.1\times10^6$. A null scan of the full location band below
$g_{a\gamma\gamma}=2.335\times10^{-14}\GeV^{-1}$ would reject the
all-dark-matter benchmark combination ($v_S{=}M_I$, $y_s{=}\sqrt2$), not
the anomaly theorem itself.

\emph{Proton decay (prior-region statements).} Under the diagnostic
priors of Table~\ref{tab:priors}, the median lifetime is
$1.1\times10^{35}$~yr; 35 per cent of sampled volume lies below the
Super-K cutoff $2.4\times10^{34}$~yr \cite{SK} and a further 14 per
cent below the Hyper-K reach $1\times10^{35}$~yr \cite{HKsnowmass};
a Hyper-K null would exclude only that additional subregion.

\FloatBarrier

\begin{table}[ht]
\centering
\caption{Diagnostic priors of the proton Monte Carlo (seed 20260801,
$10^5$ trials; spectator shift applied per trial with $y_s=\sqrt2$).}
\label{tab:priors}
\begin{tabular}{ll}
\toprule
$\alpha_s(M_Z)$ & $\mathcal{N}(0.1179,\,0.0009)$ \\
$\alpha_1^{-1},\ \alpha_2^{-1}$ & $\mathcal{N}(59.02,0.02)$,
  $\mathcal{N}(29.57,0.02)$ \\
threshold & $\Delta_4=\delta\sim\mathcal{U}(-1,1)$,\
  $\Delta_{2R}=-\tfrac23\delta$,\ $\Delta_{2L}=0$ \\
$M_X/M_{\rm GUT}$ & $\mathrm{LogU}(0.5,\,2)$ \\
$A_R$ & $\mathcal{U}(2.0,\,3.2)$ \\
$W_0$ [GeV$^2$] & $\mathcal{N}(0.1095,\,0.011)$ \\
\bottomrule
\end{tabular}
\end{table}

\section{Novelty and conclusions}
The defensible contribution now has five layers: (i) the charge-independent
five-pair lower bound, plus residue uniqueness in the identical-pair ansatz,
for this particular $\Zs$ extension; (ii) the discrete-EFT $P=12$
multi-spurion lower bound and its
explicit Spin(10) graph; (iii) the anomaly-free $U(1)_X\to\Zs$ field-theory
embedding of Section~\ref{sec:uv}; (iv) the one-pair no-go and minimal
three-pair decay-safe completion; and (v) the heavy-inclusive $P=8$
closure audit plus its repeated-pole two-loop integral in
Section~\ref{sec:matching}. These are narrow model-building and matching
results within the established discrete-gauge axion literature
\cite{BGW,DPT,DiLuzio,BDM}, not a derivation of the full ERT parameter
point from one scale and not experimental evidence for an axion.

The analytic charge ledger for the continuous embedding passes its stated
internal tests: all local anomalies
vanish, renormalizable PQ remains accidental, the axion survives the gauge
projection, and the explicitly evaluated unit-coefficient quality terms
are far below the neutron-EDM requirement. Every anomalon pair now has a
renormalizable ordinary-matter portal, and a generic mass matrix leaves
three chiral families. The completion also changes the EFT interpretation:
the old two-loop graph is $P=16$, while the decay-enabled heavy matter opens
a $P=8$ two-loop closure. Its exact benchmark shift is nevertheless four
orders below the direct dimension-21 scalar term, which dominates the
calculated terms. Unknown Wilson/flavour/Yukawa tensors remain parameters,
so ``full amplitude'' here means the complete finite momentum kernel and
mass dependence for the stated graph, not a coefficient-free prediction.
It is not an external implementation certificate. The repository model file
now uses native non-supersymmetric SARAH syntax, includes the manuscript's
gauged $U(1)_X$, and passes static catalogue, charge, Lagrangian, filter, and
hash-bound input-manifest checks. A Mathematica/SARAH process-log attestation
is not available in the present environment, so the authoritative whole-model
gate remains blocked. This also does not mean
a completed scalar-vacuum proof. The exact-$X$
renormalizable scalar
census contains 44 tensor directions and 51 real parameters. Its dense
derivative audit passes. An exact Gaussian-integer projector identity proves
that the selected-vacuum 54 and $\overline{1050}$ projector-gradient columns
vanish, and a separate exact integer/rational calculation proves the mixed
$\Phi$--$\overline{126}$ 210 column also vanishes. An exact nonzero
$13\times13$ compiler-bound minor and an exact full-row factorization prove
stationarity rank 13 and nullity 38. Normalized floating-point SVD agrees but is
used only as a diagnostic; an exact 38-vector parameter-nullspace basis is not
yet part of the G3 solver contract. Independently, an exact Gaussian-integer
tangent certificate proves that the broken SO(10) plus gauged $U(1)_X$ orbit
has rank 37. Its gauge quotient is therefore 449-dimensional and includes the
physical axion. The independent global-PQ tangent raises the full symmetry
orbit rank to 38; quotienting that flat orbit for Hessian positivity leaves a
448-dimensional massive/transverse space. PQ is not gauge-eaten. The neutral
$\Phi_{210}$
projector $P_{24}$ is exactly symmetric, idempotent, and rank 24. An exact
normalized stationary witness with coefficients $(10,1,-1/4)$ has vanishing
dense gradient and $\operatorname{tr}(P_{24}H)=288$. The former normalized-SVD
stationary constraints reject this witness, so the associated finite-cut,
common-kernel, block-SDP, and negative-trace-LP results are invalidated. The
physical $H_6$ radial coordinate is also not an exact flat: writing its
canonical real coordinate as $t=\sqrt2 h$, the exact stationary combination
$c_{O06}=-t^2$, $c_{O36,1}=10$ has vanishing first derivative and curvature
$2t^2=4h^2>0$. Adding the independent normalized $P_{24}$ witness preserves
this result and the coupling bound.

Using the corrected 11-plus-2 stationary constraints directly on the raw
orthonormal 448-dimensional massive/transverse quotient, an opt-in float64
recomputation gives common-Gram rank 448 and nullity zero. The former apparent
135-dimensional common flat subspace is reproduced only after a
reference-derived diagonal congruence with condition ratio approximately
$1.20\times10^8$ and is invalidated as a conditioning artifact. This numerical
full-rank result is not a Hessian-positivity certificate.

The corrected search nevertheless produces a new constructive frontier. A
sparse exact-$X$ coefficient vector uses 27 of the 51 real parameters, obeys
$\lVert c\rVert_\infty=73/8<4\pi$, and has $J_0=-21/200$. It is therefore
outside the former $J_0=+1$ slice, showing that this anchor was not without
loss of generality. The vector has a manifest sum-of-squares structure. In
particular, a direct Gaussian-integer tensor calculation together with exact
rational Casimir projectors proves the six-channel recoupling
\begin{equation}
 \lVert \mathcal M(\Phi)\Sigma\rVert^2 =
 40 I_1+72 I_{45}+28 I_{210}-8 I_{770}-12 I_{5940}+12 I_{8910}.
\end{equation}
A nonsingular rational six-witness evaluation fixes these coefficients
uniquely. A separate source-bound expansion of all 27 nonzero couplings proves
that the complete potential is a sum of nonnegative squares. It also proves
exact stationarity of the selected vacuum, including all 252 real
$\overline{126}$ coordinates. Thus boundedness below and selected-vacuum
stationarity are exact statements, not numerical scans.

For the $\Phi_{210}+\Delta_R$ Hessian, direct Gaussian-integer and rational
tensor contraction followed by exact $\mathbb Q(\sqrt{2})$ component
arithmetic gives $\operatorname{rank}K=278$,
$\operatorname{rank}H_\Phi=186$, and
$\operatorname{rank}(H_\Phi+K)=429$ with nullity 33. The proof assembles every
entry from integer/Gaussian-integer tensors and exact projector fractions; the
float64 live-source comparison is only a regression check. The exact
$P+\Delta_R$ gauge-orbit matrix has rank 33. After quotienting that kernel, an
explicit $26\times24$ integer Jacobian for the nonnegative $H/S/\Phi_{17}$
constraints has exact rank 19. The five remaining added directions combine
with the 33 core gauge tangents to give precisely the 38 SO(10), gauged-
$U(1)_X$, and global-PQ symmetry tangents. Hence the full Hessian has rank 448
and is positive on all transverse directions.

Combining exact stationarity, complete-potential BFB, and transverse Hessian
positivity proves that the selected symmetry orbit is a strict local minimum.
The final exact global-gap test also gives a decisive negative result for
globality.  In canonical $overline{126}$ coordinates there is a Gaussian-
integer field direction with normalized projector fractions
$(I_{54},I_{1050},I_{2772},I_{4125})=(0,0,1/2,1/2)$.  It annihilates both mixed
squares at $P$ and, after choosing $N_\Sigma=(100/99)r^2$, has energy below the selected orbit by
$25r^4/19008>0$.  The selected orbit is therefore not the global vacuum and
the present 27-parameter candidate is rejected for G3.  G3 closure remains
false, but this does not exclude the theory: the fixed-$P$ branch obeys the
exact gap--curvature relation $Delta V=-m_perp^2/8$, and its lower replacement
has the wrong stabilizer. A distinct $p:a:\omega=1:1:1$ SU(5)-singlet branch
has an explicit global $Phi+\Sigma$ SOS, the SM stabilizer, and exact
rank/nullity $429/33$.  With the neutral chiral vector
$H=(e_6+i e_7)/\sqrt2$, its 28-parameter full-field extension is exactly
stationary and BFB, has symmetry ranks $36/37/38$, and has a live positive
448-dimensional quotient with minimum eigenvalue $0.00484459$.  A separate
source-bound rational-lattice calculation and exact block LDL certificate now
prove full Hessian rank/nullity $448/38$, positive semidefiniteness, and that
the kernel is exactly the 38-dimensional symmetry orbit.  The literal
single-orbit four-form lemma is false: $F$ and $-F$ obey both projector
equations but are separated by the cubic invariant
$\operatorname{Tr}(A_\Phi^3)=\pm60$.  Exact mixed rank $252/252$ excludes the
coupled $-F$ equality branch.  The fail-closed final G3 test still requires a
global argument beyond the already certified strata.  An exact
implicit-function and equivariance argument proves that
both signed orbits are isolated local components of the full projector-zero
set.  The complete 16-real-dimensional SU(3)-fixed subspace is also classified
exactly and contains no additional orbit.  At fixed $\Phi=F$, a separate exact
off-kernel estimate proves the full $\beta=1/20$ gap for arbitrary $H$ and
  $\Sigma$, with equality only on the selected SU(5) flag orbit.  The strongest
  all-zero-residual route can also be removed exactly.  On the pure-$\Delta_R$,
  maximally negative-current stratum the live integral affine system has
  rank/nullity $168/42$ and its integer kernel splits as $35+7$.  Exact square
  completion gives a $\Phi$ radial-plus-$I_{54}$ lower bound $1/141$ and a worst
  radial/current value $-7001/995000$, hence the strict positive margin
\begin{equation}
 \frac{1}{141}-\frac{7001}{995000}
 =\frac{7859}{140295000}>0.
\end{equation}
This zero-residual result is strengthened by an exact source-bound certificate
that retains all shifted $\Phi$--$\Sigma$ and chiral $\Phi$--$H$ residuals and
covers every real $\Phi$ and all nonnegative radial variables on the same
maximally negative pure-$\Delta_R$ sector.  A single exact 4125-projector
covariant, a rational Schur-complement bound, and a three-region radial
completion prove the sharp restricted gap
\begin{equation}
 \Delta V_{\mathrm{pure}\text{-}\Delta_R}\geq \frac{1}{5000},
\end{equation}
with equality at $u=1$, $v=0$ for a signed K\"ahler-square representative.
At fixed $H=h_-$ and one explicit normalized decomposable rank-one $\Sigma$
endpoint, a separate
source-bound rational Gram/LDL certificate proves an angular anchor at least
$3/200$ and the same restricted gap $1/5000$ for all nonnegative radial
variables on a four-real-dimensional $\Phi$ sub-slice of the
16-real-dimensional SU(3)-fixed space.  This statement is only a sub-slice
theorem: it does not cover arbitrary $\Phi$ at that rank-one $\Sigma$
endpoint, the ambient SU(3)-fixed space, arbitrary $\Sigma$ orientations, or G3.
At the same fixed $H=h_-$, $\Sigma=q/4$ endpoint, a separate exact tangent
calculation identifies the common continuous stabilizer as SU(4): the joint
kernel has dimension 15, the displayed integral Lie brackets close, and all
15 skew actions on $\Phi_{210}$ are source-bound. Exact intertwiners decompose
the complexified 210 into 25 carriers, and deterministic common lowering words
align them into an exact rank-210 basis with physical conjugation and real-form
maps. An exact $5952\times551$ constraint system has rational rank 506 and
nullity 45; its displayed integral nullspace gives an explicit complete basis
of 45 real symmetric quadratic forms that commute with all 15 live actions.
Thus $\dim\!\left(\operatorname{Sym}^2(\Phi_{210})\right)^{\mathrm{SU}(4)}=45$.
The augmented homogeneous representation census is now exact as well. Its
real dimension is 22366, with 35 complex isotypic types spanning 824
irreducible copies and 22 real Schur blocks: 9 real-symmetric and 13
complex-Hermitian.  Their multiplicity matrices contain 19594 real Schur
parameters, graded $(1,4,90,1414,18085)$, while the invariant quartic target
has 6585 rows, graded $(1,4,45,478,6057)$. A universal grade-preserving
rational polarization section proves the abstract invariant multiplication
map surjective, with total rank 6585 and kernel dimension 13009.
This is an abstract census, not the complete Gram problem. The 35 aligned
isotypic embeddings spanning all 824 copies, ordered invariant cubic and
quartic coordinates, physical G3 gap target, the $6585\times19594$ coordinate
Schur matrix, and PSD feasibility or infeasibility certificates remain to be
constructed. No arbitrary-$\Phi$ bound, G3 closure, or whole-model conclusion
follows from the abstract surjectivity statement.
Consequently G3 remains open only because this uniform coercive estimate has
not yet been extended to arbitrary non-pure-$\Delta_R$ $\Sigma$ orientations. The
historical finite-cut and SDP solvers remain non-certifying and are not used in
the local-minimum proof. Therefore this work neither validates nor excludes
the complete model.

The v16--v20 audit also records negative results: the v15 dimension-9
one-insertion potential and its $\mathbb Z_3$ cure were incorrect, and the
alternate v18 scalar-only UV-quality proof is incomplete. V19's exact
stable-anomalon caveat is resolved rather than hidden: its charge set is
replaced. The remaining scientific tasks are an independent
representation-theory and diagrammatic review, broken-phase Clebsch and
flavour fits, RG evolution of the Wilson tensors, a radiatively stable
$v_\Phi/v_S$ hierarchy, a thermal-restoration calculation, an ultimate
quantum-gravity origin, and the nonstandard string-network calculation.
The release retains the v17 65-check and v19 59-check regression engines
and adds a v20 42-check engine with a controlled nonzero-exit path and
machine-readable verdict.

\appendix
\section{Reproduction formulas}
The rephasing diagnostic referenced in Section~\ref{sec:quality} is
\begin{equation}
M =
\begin{pmatrix}
0 & 2 & 1 & 0 & 0 & 0\\
0 & 2 & 0 & 1 & 0 & 0\\
1 & 0 & 0 & 0 & 1 & 1\\
2 & 0 & 2 & 2 & 0 & 0\\
-3 & 1 & 0 & 0 & 0 & 1\\
4 & 1 & 1 & 0 & 1 & 0\\
-2 & 1 & 0 & 0 & 1 & 2\\
-2 & 0 & 0 & 0 & 4 & 0
\end{pmatrix}
\end{equation}
over $(X_S, X_F, X_{10}, X_{\overline{126}}, X_a, X_b)$, rows: the two
Yukawas, spectator mass, ERT lock, $(S^\dagger)^3 16_F\overline{16}_s$,
$S^4 16_F 16_s 10_H$,
$(S^\dagger)^2(16_F\overline{16}_s)(16_s\overline{16}_s)$,
$(S^\dagger)^2(16_s16_s)_{10}(16_s16_s)_{10}$. The proton width used in
the Monte Carlo is
\begin{equation}
\Gamma(p\to e^+\pi^0) = \frac{m_p}{32\pi}
\left[1-\frac{m_\pi^2}{m_p^2}\right]^2
\left(\frac{4\pi\alpha_{\rm GUT}}{M_X^2}\right)^{\!2}
A_R^2\,|W_0|^2\,\bigl[1+(1+|V_{ud}|^2)^2\bigr],
\end{equation}
with $W_0$ from lattice QCD \cite{Yoo} and $\tau_p = \hbar/\Gamma$.

\clearpage

\begin{thebibliography}{99}
\small
\sloppy
\bibitem{ERT} A.~Ernst, A.~Ringwald and C.~Tamarit,
  JHEP \textbf{02} (2018) 103, arXiv:1801.04906.
\bibitem{Hsieh} C.-T.~Hsieh, arXiv:1808.02881.
\bibitem{KW} L.~M.~Krauss and F.~Wilczek,
  Phys.\ Rev.\ Lett.\ \textbf{62} (1989) 1221.
\bibitem{BanksDine} T.~Banks and M.~Dine,
  Phys.\ Rev.\ D \textbf{45} (1992) 1424, arXiv:hep-th/9109045.
\bibitem{Ibanez93} L.~E.~Ib\'a\~nez,
  Nucl.\ Phys.\ B \textbf{398} (1993) 301, arXiv:hep-ph/9210211.
\bibitem{BGW} K.~S.~Babu, I.~Gogoladze and K.~Wang,
  Phys.\ Lett.\ B \textbf{560} (2003) 214, arXiv:hep-ph/0212339.
\bibitem{DPT} A.~G.~Dias, V.~Pleitez and M.~D.~Tonasse,
  Phys.\ Rev.\ D \textbf{67} (2003) 095008, arXiv:hep-ph/0211107.
\bibitem{DiLuzio} L.~Di Luzio, JHEP \textbf{11} (2020) 074,
  arXiv:2008.09119.
\bibitem{BDM} K.~S.~Babu, B.~Dutta and R.~N.~Mohapatra,
  Phys.\ Rev.\ Lett.\ \textbf{134} (2025) 111803, arXiv:2410.07323.
\bibitem{BDMquality} K.~S.~Babu, B.~Dutta and R.~N.~Mohapatra,
  ``Accidental Peccei--Quinn Symmetry From Gauged U(1) and a High
  Quality Axion,'' arXiv:2412.21157.
\bibitem{Bonnefoy} Q.~Bonnefoy, Phys.\ Rev.\ D \textbf{108} (2023)
  035023, arXiv:2212.00102.
\bibitem{Djouadi} A.~Djouadi, R.~Fonseca, R.~Ouyang and M.~Raidal,
  Eur.\ Phys.\ J.\ C \textbf{83} (2023) 654, arXiv:2212.11315.
\bibitem{Chang} D.~Chang, T.~Fukuyama, Y.-Y.~Keum, T.~Kikuchi and
  N.~Okada, ``Perturbative SO(10) Grand Unification,''
  arXiv:hep-ph/0412011.
\bibitem{LRS} Q.~Lu, M.~Reece and Z.~Sun, JHEP \textbf{07} (2024) 227,
  arXiv:2312.07650.
\bibitem{KKT} T.~Kobayashi, R.~Kurematsu and F.~Takahashi,
  JCAP \textbf{09} (2013) 032, arXiv:1304.0922.
\bibitem{GR25} P.~W.~Graham and D.~Racco, JHEP \textbf{12} (2025) 028,
  arXiv:2506.03348.
\bibitem{Benabou} J.~Benabou, M.~Buschmann, J.~W.~Foster and
  B.~R.~Safdi, arXiv:2412.08699.
\bibitem{Saikawa} K.~Saikawa, J.~Redondo, A.~Vaquero and
  M.~Kaltschmidt, ``Spectrum of global string networks and the axion
  dark matter mass,'' JCAP \textbf{10} (2024) 043, arXiv:2401.17253.
\bibitem{Borsanyi} S.~Bors\'anyi et al., ``Calculation of the axion
  mass based on high-temperature lattice quantum chromodynamics,''
  Nature \textbf{539} (2016) 69, arXiv:1606.07494.
\bibitem{GrillidiCortona} G.~Grilli di Cortona et al.,
  JHEP \textbf{01} (2016) 034, arXiv:1511.02867.
\bibitem{Yoo} J.-S.~Yoo et al., Phys.\ Rev.\ D \textbf{105} (2022)
  074501, arXiv:2111.01608.
\bibitem{DALI} DALI Collaboration, arXiv:2603.21951 (2026).
\bibitem{DALIcomment} Comment on the DALI sensitivity,
  arXiv:2607.28095 (2026).
\bibitem{IR} L.~E.~Ib\'a\~nez and G.~G.~Ross, ``Discrete gauge
  symmetries and the origin of baryon and lepton number conservation,''
  Nucl.\ Phys.\ B \textbf{368} (1992) 3 (mixed-gauge discrete
  anomalies).
\bibitem{CAST} CAST Collaboration, Phys.\ Rev.\ Lett.\ \textbf{133}
  (2024) 221005, arXiv:2406.16840.
\bibitem{MADMAX} P.~Brun et al.\ (MADMAX), Eur.\ Phys.\ J.\ C
  \textbf{79} (2019) 186.
\bibitem{ALPHA} A.~J.~Millar et al., Phys.\ Rev.\ D \textbf{107}
  (2023) 055013, arXiv:2210.00017.
\bibitem{ORGAN} B.~T.~McAllister et al., Phys.\ Dark Univ.\ \textbf{18}
  (2017) 67, arXiv:1706.00209.
\bibitem{SK} A.~Takenaka et al.\ (Super-Kamiokande),
  Phys.\ Rev.\ D \textbf{102} (2020) 112011.
\bibitem{HKsnowmass} J.~Bian et al.\ (Hyper-Kamiokande),
  arXiv:2203.02029.
\end{thebibliography}

\end{document}
